\section{Transición tecnológica y Wedge: por qué Cloud Identity}

McKinnon entiende la oportunidad de emprender dentro de la cloud transition: cuando las aplicaciones empresariales pasan de centros de datos propios a SaaS, el antiguo perímetro de red y los directorios locales ya no cubren de forma natural todos los recursos. Una empresa nueva no puede replicar de frente el mercado antiguo en el que el incumbent es más fuerte; debe encontrar el nuevo control point que la transición deja expuesto y entrar con un wedge por el que el cliente esté dispuesto a pagar de inmediato.

\subsection{Del Monitoring Pivot a SSO}

El equipo empezó con cloud application monitoring, pero el buyer feedback fue tibio; en cambio, ante varias cuentas de SaaS como Gmail y Dropbox, los inicios de sesión repetidos, las contraseñas y los permisos de quienes dejan la empresa se volvieron un problema inmediato para los empleados. \term{single sign-on} (inicio de sesión único, SSO) permite al usuario acceder a varias aplicaciones a través de un IdP de confianza, lo que reduce las credenciales repetidas y centraliza la policy. SSO es solo un wedge, pero por su naturaleza exige conectarse con muchas aplicaciones, directorios y protocolos.

\lecturefigure{01-transition-window.png}{La transición tecnológica crea nuevos puntos de control; el identity wedge crece desde el problema inmediato del inicio de sesión hasta convertirse en infraestructura compartida.}{Subtítulos locales 00:00--05:20; redibujo conceptual.}

\begin{knowledgebox}{Lectura de la figura: Transition, pain y platform deben cumplirse a la vez}
La cloud transition cambia dónde están las aplicaciones y cómo se compran; la credential fragmentation genera un pain concreto; SSO ofrece un resultado de corto plazo por el que se paga; directory, policy y lifecycle dan al producto profundidad de plataforma. Con transition pero sin buyer pain, se convierte en una investigación de largo plazo; con un problema pequeño pero sin superficie de expansión, se queda en feature; con platform vision pero sin wedge, es difícil obtener el primer tráfico de producción.
\end{knowledgebox}

\teachervoice{McKinnon no idealiza la primera idea: una vez construido el monitoring prototype, los clientes potenciales no mostraron entusiasmo, y por eso el equipo giró hacia «hacer más simple el inicio de sesión en SaaS». La nota de clase: primero buscar una willingness to pay real y después demostrar que detrás de la pequeña entrada existe una infrastructure con suficiente profundidad.}

\subsection{El iceberg de la infraestructura: la complejidad detrás de una experiencia simple}

Que un empleado haga clic en el ícono de una aplicación parece solo «escribir la contraseña una vez menos»; sin embargo, el sistema debe sincronizar el estado del usuario, elegir la fuente de identidad, aplicar las políticas de autenticación, emitir la assertion/token, establecer la session, gestionar logout y revoke, y actualizar los entitlements cuando la organización cambia. El efecto acumulativo del identity wedge proviene de que estos primitive compartidos pueden atender cada vez más aplicaciones, no de la página de inicio de sesión en sí.

\begin{importantbox}{Las dos restricciones del Wedge}
Una entrada de emprendimiento en infraestructura necesita a la vez: \textbf{today value}, que el cliente esté dispuesto a pagar hoy por menos fricción al iniciar sesión, un onboarding más rápido o un offboarding más controlable; y \textbf{tomorrow depth}, que la entrada pueda extenderse a directory, MFA, policy, provisioning, governance y threat response. Si falta cualquiera de los dos, el producto se queda sin flujo de efectivo o sin espacio de plataforma a largo plazo.
\end{importantbox}

\subsection{Resumen de la sección}

La oportunidad de cloud identity proviene del cambio en los límites de las aplicaciones. SSO es la puerta de entrada que se puede vender; identity lifecycle y policy son la plataforma posterior. El Pivot se fundamenta en la evidencia de los clientes, no en el costo hundido del fundador en su primer prototype.
